\notepart{The one-dimensional matrix product state approximation}{part:mps}

\section{Statement, and the two-line route}\label{sec:shortcut}

\begin{inwords}
Theorem~\ref{thm:B} asks for a single global matrix product state approximating
the ground vector, rather than an entropy bound on each cut individually.
An entropy bound alone does not yield a rank truncation because von Neumann
entropy controls only the mean of the entanglement spectrum, not the weight
of the small Schmidt eigenvalues discarded when truncating a cut; in two
dimensions, moreover, truncations across different cuts do not assemble into
a consistent tensor network, and a distributed protocol is needed to bridge
the gap. In one dimension there is a classical shortcut: if the Schmidt
spectrum across every prefix cut has a rapidly decaying tail, then
truncating the cuts sequentially from one end of the chain to the other
produces a matrix product state whose global vector error is bounded by the
sum of the square roots of the discarded tail weights (Verstraete--Cirac).
The marginal tail bound of Section~\ref{sec:tail} provides precisely such an
exponential tail at a rate $u_0$ depending only on $q,R,J/\Delta$, while the
area law fixes the mean of the spectrum. Combining the two yields
Theorem~\ref{thm:B} in a few lines, with bond dimension $CL^{4/u_0}$.

We record this sequential-truncation proof first because it is direct and
elementary. The remainder of Part~\ref{part:mps} then develops the
one-dimensional \emph{protocol} proof, which avoids sequential cut
truncations altogether and is the construction that generalises to two
dimensions.
\end{inwords}

\precisely
\begin{figure*}[t]
\centering
\begin{adjustbox}{max width=\linewidth}
\(\displaystyle \Phi\;=\;
  \begin{tenkz}[size=l, rows={wire,wire}, cols=6, bonds=none]
    \tn[at=(2,1), name=a1, skin=box, ports={0:virtual, 90:physical}]{A^{[1]}}
    \tn[at=(2,2), name=a2, skin=box, ports={180:virtual, 0:virtual, 90:physical}]{A^{[2]}}
    \tn[at=(2,3), name=a3, skin=box, ports={180:virtual, 0:virtual, 90:physical}]{A^{[3]}}
    \tn[at=(2,4), name=a4, skin=box, ports={180:virtual, 0:virtual, 90:physical}]{A^{[4]}}
    \tn[at=(2,5), name=a5, skin=box, ports={180:virtual, 0:virtual, 90:physical}]{A^{[5]}}
    \tn[at=(2,6), name=a6, skin=box, ports={180:virtual, 90:physical}]{A^{[6]}}
    \tn[at=(1,1), name=p1, skin=none, ports={270:physical, 90:physical}]{}
    \tn[at=(1,2), name=p2, skin=none, ports={270:physical, 90:physical}]{}
    \tn[at=(1,3), name=p3, skin=none, ports={270:physical, 90:physical:$q$}]{}
    \tn[at=(1,4), name=p4, skin=none, ports={270:physical, 90:physical}]{}
    \tn[at=(1,5), name=p5, skin=none, ports={270:physical, 90:physical}]{}
    \tn[at=(1,6), name=p6, skin=none, ports={270:physical, 90:physical}]{}
    \tnwire{a1.90}{p1.270} \tnwire{a2.90}{p2.270} \tnwire{a3.90}{p3.270}
    \tnwire{a4.90}{p4.270} \tnwire{a5.90}{p5.270} \tnwire{a6.90}{p6.270}
    \tnbond{a1.0}{a2.180} \tnbond{a2.0}{a3.180} \tnbond[name=d]{a3.0}{a4.180}
    \tnbond{a4.0}{a5.180} \tnbond{a5.0}{a6.180}
    \tnmark[form=label, label pos=270]{on d 0.5}{$D$}
  \end{tenkz}\)

\end{adjustbox}
\caption{A matrix product state with open boundary conditions: one tensor per site, bonds of dimension $D$ between neighbours, one physical leg per site.}\label{fig:mps-chain}
\end{figure*}

\begin{proposition}[Sequential truncation]\label{prop:VC}
Let $\Om$ be a unit vector on $\cH_\Lam$, $\Lam=\{1,\dots,L\}$. For
$1\le k\le L-1$ let $\Pi_k$ be the projection of $\cH_{[1,k]}$ onto the
span of the $\min\{D,\dim\cH_{[1,k]}\}$ eigenvectors of $\rho_{[1,k]}$ with
largest eigenvalues (so $\Pi_k=\id$ when $D\ge\dim\cH_{[1,k]}$), let $\epsilon_k=\langle\Om,(\id-\Pi_k\otimes\id)\Om\rangle$ be the discarded
mass, and put
\[
 \Phi=(\Pi_1\otimes\id)(\Pi_2\otimes\id)\cdots(\Pi_{L-1}\otimes\id)\,\Om .
\]
Then $\Phi$ is a matrix product state with open boundary conditions and
bond dimension at most $D$, and $\norm{\Phi-\Om}\le\sum_k\sqrt{\epsilon_k}$.
\end{proposition}
\begin{proof}
Write $P_k=\Pi_k\otimes\id$. Telescoping gives
$\Phi-\Om=\sum_{k=1}^{L-1}P_1\cdots P_{k-1}(P_k-\id)\Om$, and since each
prefix $P_1\cdots P_{k-1}$ is a contraction, the triangle inequality gives
$\norm{\Phi-\Om}\le\sum_k\norm{(P_k-\id)\Om}=\sum_k\sqrt{\epsilon_k}$. To
bound the bond dimension, fix a cut $k$. The intermediate vector
$P_kP_{k+1}\cdots P_{L-1}\Om$ lies in
$\operatorname{ran}\Pi_k\otimes\cH_{[k+1,L]}$ and therefore has Schmidt rank
at most $D$ across cut $k$; the remaining factors $P_1,\dots,P_{k-1}$ act
solely on $[1,k-1]$ to the left of cut $k$ and cannot increase that Schmidt
rank. A vector with Schmidt rank at most $D$ across every cut is a matrix
product state of bond dimension at most $D$.
\end{proof}

\begin{figure*}[t]
\centering
\begin{adjustbox}{max width=\linewidth}
\(\displaystyle \Phi\;=\;
  \begin{tenkz}[rows={wire,wire,wire,wire,wire}, cols=5, bonds=none]
    \tn[at=(5,1), skin=box, wide=5, name=o,
        ports={90@1:physical, 90@2:physical, 90@3:physical, 90@4:physical, 90@5:physical}]{\Omega}
    \tn[at=(4,1), skin=box, wide=4, name=p4,
        ports={270@1:physical, 270@2:physical, 270@3:physical, 270@4:physical,
               90@1:physical, 90@2:physical, 90@3:physical, 90@4:physical}]{P_4}
    \tn[at=(3,1), skin=box, wide=3, name=p3,
        ports={270@1:physical, 270@2:physical, 270@3:physical,
               90@1:physical, 90@2:physical, 90@3:physical}]{P_3}
    \tn[at=(2,1), skin=box, wide=2, name=p2,
        ports={270@1:physical, 270@2:physical, 90@1:physical, 90@2:physical}]{P_2}
    \tn[at=(1,1), skin=box, name=p1,
        ports={270:physical, 90:physical:$1$}]{P_1}
    \tn[at=(1,2), skin=none, name=t2, ports={270:physical, 90:physical:$2$}]{}
    \tn[at=(1,3), skin=none, name=t3, ports={270:physical, 90:physical:$3$}]{}
    \tn[at=(1,4), skin=none, name=t4, ports={270:physical, 90:physical:$4$}]{}
    \tn[at=(1,5), skin=none, name=t5, ports={270:physical, 90:physical:$5$}]{}
    \tnwire{o.90@1}{p4.270@1} \tnwire{o.90@2}{p4.270@2}
    \tnwire{o.90@3}{p4.270@3} \tnwire{o.90@4}{p4.270@4}
    \tnwire{p4.90@1}{p3.270@1} \tnwire{p4.90@2}{p3.270@2} \tnwire{p4.90@3}{p3.270@3}
    \tnwire{p3.90@1}{p2.270@1} \tnwire{p3.90@2}{p2.270@2}
    \tnwire{p2.90@1}{p1.270}
    \tnwire{p2.90@2}{t2.270}
    \tnwire{p3.90@3}{t3.270}
    \tnwire{p4.90@4}{t4.270}
    \tnwire{o.90@5}{t5.270}
  \end{tenkz}\)

\end{adjustbox}
\caption{Sequential truncation, drawn for $L=5$. The projector $P_k=\Pi_k\otimes 1$ onto the top $D$ Schmidt vectors of the prefix $[1,k]$ is applied for $k=L-1$ first and $k=1$ last; the output of $P_k$ lies in $\mathrm{ran}\,\Pi_k\otimes\mathcal H_{[k+1,L]}$, so across every cut the result has Schmidt rank at most $D$.}\label{fig:truncation-sequence}
\end{figure*}

\begin{theorem}[Theorem~\ref{thm:B}, two-line route]\label{thm:Bshort}
Under the hypotheses of Theorem~\ref{thm:B}, for every $L\ge2$ there is
a matrix product state $\Phi$ with bond dimension at most
$2\,\e^{C_{\mathrm{int}}}\,(32\e^{\e/2}L^4)^{1/u_0}$ and
$\min_\theta\norm{\Phi/\norm\Phi-\e^{i\theta}\Om}\le L^{-1}$, where $u_0$ is
\eqref{eq:theta} and $C_{\mathrm{int}}$ is \eqref{eq:final}.
\end{theorem}
\begin{proof}
Fix a cut $k$ and put $X=[1,k]$. By Lemma~\ref{lem:tail}, the eigenvalues
$p$ of $\rho_X$ with $-\log p>S(X)+w$ have total mass at most
$2\e^{\e/2}\e^{-u_0w}$, while the eigenvalues with $-\log p\le S(X)+w$ each
satisfy $p\ge\e^{-S(X)-w}$ and therefore number at most $\e^{S(X)+w}$.
Choose $w$ so that $2\e^{\e/2}\e^{-u_0w}=(16L^4)^{-1}$ and set
$D=\lceil\e^{C_{\mathrm{int}}+w}\rceil\ge\e^{S(X)+w}$ by Theorem~\ref{thm:A}.
Then $\epsilon_k\le(16L^4)^{-1}$ for every $k$, and Proposition~\ref{prop:VC}
gives $\norm{\Phi-\Om}\le(L-1)/(4L^2)\le1/(4L)$. In particular $\Phi\ne0$
and $\norm{\Phi/\norm\Phi-\Om}\le2\norm{\Phi-\Om}\le L^{-1}$. The bond
dimension satisfies $D\le2\e^{C_{\mathrm{int}}+w}$ with
$\e^{w}=(32\e^{\e/2}L^4)^{1/u_0}$.
\end{proof}

\begin{remark}
The exponent $4/u_0$ is of order $\sqrt{(1+J/\Delta)\cB_R}$. Whereas
Hastings' original construction yields a polynomial bond dimension whose
exponent is exponential in $J/\Delta$, the marginal tail bound gives a
square-root dependence on $J/\Delta$. This exponent is a feature of the
direct sequential-truncation route rather than of the protocol below.
\end{remark}

\subsection*{Plan of the protocol route}

\begin{inwords}
The distributed protocol proof has four layers:
\begin{enumerate}[leftmargin=2em]
\item \emph{Information at logarithmic separation.} Two regions of the
chain separated by $D\log L$ sites share at most $L^{-p}$ nats of mutual
information. This is where the protocol route uses the
spectral gap. In two dimensions, establishing logarithmic
separation requires a coarse collar bound, a cell conditional-information
lemma, and an angular induction; in one dimension, it follows directly from
the area law and a single two-replica information-reset lemma
(Section~\ref{sec:logsep}).
\item \emph{Ownership moves.} Three elementary operations redistribute site
registers among parties while keeping the global state close to a reference
state: a \emph{birth} (handing an interior interval to a new owner), its
reverse, a \emph{death}, and an \emph{exchange} of ownership assignments
between two copies of the ground state inside a window. Each move is a
contraction built from private maps and at most one normalised shared
bipartite vector or covector, and each incurs $O(L^{-p/2})$ vector error
whenever the moved region has mutual information at most $L^{-p}$ with all
third-party registers (Section~\ref{sec:moves}).
\item \emph{The sweep schedule.} A left-to-right sweep with a moving
remainder party creates the logarithmic block owners one by one. The only
obstruction is dissolving a narrow remainder band wedged between two
distinct neighbouring owners; this is accomplished by a fixed sequence of
moves using an auxiliary copy of the ground state, copied from the edge
construction of the two-dimensional paper \cite{PP}, so that every party
participates in a bounded number of moves (Section~\ref{sec:sweep}).
\item \emph{Compression and assembly.} The sweep is a linear calculation
whose intermediate private registers have arbitrarily large dimension.
Replacing each shared bipartite resource in the density calculation by an
empirical average of polynomially many random product operators approximates
the output physical density in trace norm; the error bound is independent
of the private dimensions because each party participates in boundedly many
gates. Extracting a single computational-basis column of the resulting
operator network and placing its nodes along the chain yields the matrix
product state (Sections~\ref{sec:compression}, \ref{sec:routing}).
\end{enumerate}
\end{inwords}

\section{Information at logarithmic separation}\label{sec:logsep}

\subsection{Purifications, fidelity, and splitting}

\begin{inwords}
To change the ownership of a region $T$ without disturbing the global state,
we must decouple $T$ from the registers $E$ held by third parties using an
operation supported entirely on the separating buffer $U=\Lam\setminus(T\cup E)$.
Small mutual information $I(T:E)$ means that the reduced state $\rho_{TE}$
is close in relative entropy to the product state $\rho_T\otimes\rho_E$.
Relative entropy controls Uhlmann fidelity, which equals the maximum overlap
between purifications. Since $U$ purifies $TE$ in the pure ground state
$\Om$, Uhlmann's theorem guarantees an isometry $V$ acting on the buffer
$U$ alone that maps $\Om$ into nearly a tensor product $s\otimes s'$ of a
purification of $\rho_T$ and a purification of $\rho_E$. Any party holding
$U$ can execute this decoupling privately; this single mechanism underlies
all births, deaths, and exchanges.
\end{inwords}

\precisely
For density operators $\rho,\sigma$ put $F(\rho,\sigma)=\norm{\sqrt\rho\sqrt\sigma}_1$
and $D(\rho\|\sigma)=\Tr\rho(\log\rho-\log\sigma)$, taken to be infinite
when the support of $\rho$ is not contained in that of $\sigma$.

\begin{lemma}[Fidelity, information, Uhlmann]\label{lem:fidelity}
For normalised finite-dimensional states, $F(\rho,\sigma)\ge\e^{-D(\rho\|\sigma)/2}$,
and $D(\rho_{AC}\|\rho_A\otimes\rho_C)=I(A:C)$. If $\xi$ purifies $\rho$ on
$\cH\otimes\cK$ and $\eta$ purifies $\sigma$ on $\cH\otimes\cK'$, then after
enlarging $\cK'$ if necessary there is an isometry $V:\cK\to\cK'$ with
$\abs{\langle\eta,(\id\otimes V)\xi\rangle}=F(\rho,\sigma)$.
\end{lemma}
\begin{proof}
Let $(p_i,\ket i)$ and $(q_j,\ket{j'})$ be eigenpairs of $\rho$ and $\sigma$.
If $D(\rho\|\sigma)<\infty$, then $w_{ij}=p_i\abs{\langle i|j'\rangle}^2$ defines
a probability distribution on the pairs with $p_iq_j>0$, and
$D(\rho\|\sigma)=\sum_{ij}w_{ij}\log(p_i/q_j)$. By convexity of the exponential,
$\Tr\sqrt\rho\sqrt\sigma=\sum_{ij}w_{ij}\exp(-\tfrac12\log(p_i/q_j))\ge\e^{-D(\rho\|\sigma)/2}$,
and the trace norm dominates the trace. The identity
$D(\rho_{AC}\|\rho_A\otimes\rho_C)=I(A:C)$ follows from expanding
$\log(\rho_A\otimes\rho_C)$. The purification statement is Uhlmann's
theorem: the overlaps $\Tr(\sqrt\rho\sqrt\sigma U)$ over unitaries $U$ on a
common purifying space have maximum modulus $\norm{\sqrt\rho\sqrt\sigma}_1$
by polar decomposition, and any two purifications of the same state are
related by an isometry on the purifying factor.
\end{proof}

\begin{lemma}[Splitting]\label{lem:splitting}
Let $T,U,E$ partition $\Lam$ and $b=I_\Om(T:E)$. There are finite-dimensional
spaces $\cB_T,\cB_E$, an isometry $V:\cH_U\to\cB_T\otimes\cB_E$ and unit vectors
$s\in\cH_T\otimes\cB_T$, $s'\in\cH_E\otimes\cB_E$ with
\[
 \norm{(\id_{TE}\otimes V)\Om-s\otimes s'}\le\sqrt{2(1-\e^{-b/2})}\le\sqrt b .
\]
\end{lemma}
\begin{proof}
Because $\rho_{TE}$ is supported in $\supp\rho_T\otimes\supp\rho_E$ (the
projector onto the kernel of a marginal, tensored with the identity, has
zero expectation in the positive joint state), the relative entropy
$D(\rho_{TE}\|\rho_T\otimes\rho_E)=b$ is finite. Choose purifications $s$
of $\rho_T$ and $s'$ of $\rho_E$, enlarge their purifying spaces so that
$\cB_T\otimes\cB_E$ contains an isometric copy of $\cH_U$, and apply
Lemma~\ref{lem:fidelity}: there is an isometry $V$ on $\cH_U$ such that
$(\id_{TE}\otimes V)\Om$ has real overlap at least $\e^{-b/2}$ with
$s\otimes s'$. Two unit vectors with real overlap $F$ have Euclidean
distance $\sqrt{2(1-F)}$.
\end{proof}

\begin{figure}[t]
\centering
\begin{adjustbox}{max width=\linewidth}
\(\begin{tenkz}[size=l, rows={wire,wire}, cols=4, bonds=none]
  \tn[at=(2,1), skin=box, wide=4, name=o, ports={90@1:physical:$T$, 90@2:physical, 90@4:physical:$E$}]{\Omega}
  \tn[at=(1,2), skin=box, wide=2, name=v, ports={270@1:physical, 90@1:physical, 90@2:physical}]{V}
  \tnwire{o.90@2}{v.270@1}
\end{tenkz}
\approx
\begin{tenkz}[size=l, rows={wire}, cols=4, bonds=none]
  \tn[at=(1,1), skin=box, wide=2, ports={90@1:physical:$T$, 90@2:physical}]{s}
  \tn[at=(1,3), skin=box, wide=2, ports={90@1:physical, 90@2:physical:$E$}]{s'}
\end{tenkz}\)

\end{adjustbox}
\caption{The splitting lemma as a diagram: the isometry $V$ on the buffer $U$ turns $\Om$ into nearly a product of a purification $s$ of $\rho_T$ and a purification $s'$ of $\rho_E$.}\label{fig:split}
\end{figure}

\subsection{The information reset}\label{sec:reset}

\begin{inwords}
The area law gives $I(R:C)=O(1)$ between an interval $R$ and any disjoint
region $C$, never $o(1)$. To assemble a global tensor network from local
operations across the chain, however, errors must be suppressed to an
arbitrary inverse power $L^{-p}$ so that they remain small when summed over
$O(L)$ steps. The information-reset lemma bridges this gap by shrinking $R$
and $C$ inward to cores $R_0\subseteq R$ and $C_0\subseteq C$ separated from
their outer boundaries by a margin of width $w=O(\log L)$.

The physical mechanism is a two-replica swap. Take two identical copies of
the ground vector, $\Psi=\Om\otimes\Om$, and swap the two copies on $R$ to
obtain $\Phi=F_R\Psi$. In the first replica of $\Phi$, the core $R_0\subseteq R$
has been swapped in from the second copy of $\Om$ while $C_0\subseteq C$
remains from the first, so the two cores originate from independent copies and
have the exact product marginal $\rho_{R_0}\otimes\rho_{C_0}$. If a unitary $U$
supported entirely in the buffer region between $R_0$ and $C_0$ maps $\Psi$
to $\Phi$ up to error $O(L^{-r_*})$, then the joint marginal of the cores in
$\Psi$ agrees with $\rho_{R_0}\otimes\rho_{C_0}$ up to $O(L^{-r_*})$, which
implies $I_\Om(R_0:C_0)\le L^{-p}$.

The unitary $U$ is constructed in three steps. First, the bound $I(R:C)\le b$
and the splitting lemma produce a Hermitian contraction $W$ on the buffer
$Q=\Lam\setminus(R\cup C)$ of both copies with macroscopic overlap
$z=\langle\Phi,W\Psi\rangle\ge\e^{-2b}$. Because $W\Psi$ still carries excited
components, $W$ is not yet close to a unitary mapping $\Psi$ to $\Phi$.
Second, a two-sided Gaussian spectral filter,
$M=\int g_h(t)\e^{it\mathbb H'}W\e^{-it\mathbb H}dt$, damps out excited states
in both directions, projecting $W\Psi$ onto $z\Phi$ and $W^\dagger\Phi$ onto
$z\Psi$ with error $\e^{-h\Delta^2/2}$, because $\Psi$ and $\Phi$ are the unique
gapped ground vectors of the doubled Hamiltonians $\mathbb H$ and
$\mathbb H'=F_R\mathbb HF_R$. Third, because time evolution spreads operator
supports, a Lieb--Robinson bound localises $M$ to within distance
$\ell=O(1+b+\log L)$ of the buffer $Q$ and of the boundary zone $Z_R$ where
$\mathbb H$ and $\mathbb H'$ differ. Rescaling the localised operator makes it
almost isometric on $\Psi$, and its polar unitary $U$ remains supported strictly
outside both cores. The lemma is proved exactly as in the two-dimensional
source; we only insert the one-dimensional locality bound.
\end{inwords}

\precisely
To control the spatial spread of the Gaussian filter, we first record the standard
Lieb--Robinson propagation estimate on the doubled chain. Let
$\mathbb H$ be a Hamiltonian on $\bigotimes_{x\in\Lam}(\C^q\otimes\C^q)$
written as a sum of terms $k_Z$ supported on sets $Z$ of diameter at most
$R$ with $\sum_{Z\ni x}\norm{k_Z}\le b_0$ for every $x$, and let
$\tau_t(A)=\e^{it\mathbb H}A\e^{-it\mathbb H}$.

\begin{lemma}[Lieb--Robinson bound on a chain]\label{lem:LR}
There are $C,v,\mu>0$ depending only on $R$ and $b_0$ such that for every
operator $A$ supported in $S\subseteq\Lam$, every site $y$, and every
operator $B_y$ supported at $y$,
\[
 \norm{[\tau_t(A),B_y]}\le C\abs S\,\norm A\,\norm{B_y}\,\e^{v\abs t-\mu\,\dist(S,y)} .
\]
Consequently, if $\mathcal P_D$ denotes the conditional expectation onto
operators supported in $D$ (the average over independent Haar unitaries at
the sites outside $D$), then for every $\ell\ge0$,
\[
 \norm{\tau_t(A)-\mathcal P_{N_\ell(S)}\tau_t(A)}\le C'\abs S^2\,\norm A\,\e^{v\abs t-\mu\ell}.
\]
\end{lemma}
\begin{proof}
For $t\ge0$ put $c_y(S,t)=\sup\{\norm{[\tau_t(A),B_y]}:\supp A\subseteq S,\norm A\le1\}$.
Differentiating and using the Jacobi identity shows that $\frac{d}{dt}[\tau_t(A),B_y]$
equals $i[\sum_{Z\cap S\ne\varnothing}\tau_t(k_Z),[\tau_t(A),B_y]]$
plus a sum over $Z\cap S\ne\varnothing$ of $-i[\tau_t(A),[\tau_t(k_Z),B_y]]$. The
first term is removed by a unitary change of frame, yielding the integral inequality
$c_y(S,t)\le c_y(S,0)+2\sum_{Z\cap S\ne\varnothing}\norm{k_Z}\int_0^tc_y(Z,u)\,du$.
Iterating this inequality, and using the trivial bound $c_y(Z,u)\le2\norm{B_y}$ to see
that the remainder after $m$ iterations vanishes as $m\to\infty$, gives
\[
 c_y(S,t)\le 2\norm{B_y}\sum_{m\ge0}\frac{(2t)^m}{m!}
 \sum_{\substack{Z_1\cap S\ne\varnothing,\ Z_{i+1}\cap Z_i\ne\varnothing\\ y\in Z_m}}
 \prod_{i=1}^m\norm{k_{Z_i}},
\]
with the $m=0$ term equal to $\mathbf 1_{\{y\in S\}}$. Since the norms of
the interaction terms meeting any set of $n$ sites sum to at most $nb_0$, the inner
sum is at most $\abs S\,b_0\,((R+1)b_0)^{m-1}\le\abs S\kappa^m$ with
$\kappa=(R+1)b_0$. Moreover, a chain of $m$ supports of diameter at most $R$
connecting $S$ to $y$ requires $mR\ge\dist(S,y)$. Hence,
for every $\mu>0$,
$c_y(S,t)\le 2\abs S\norm{B_y}\e^{-\mu\dist(S,y)}\sum_m(2\kappa\e^{\mu R}t)^m/m!$,
which gives the first bound with $v=2\kappa\e^{\mu R}$; negative times are
identical. For the second
bound, expressing $\mathcal P_D$ as a product of commuting single-site Haar twirls
$T_y$ and telescoping gives
$\norm{B-\mathcal P_DB}\le\sum_{y\notin D}\sup_{U_y}\norm{[B,U_y]}$. Because at most
$2\abs S$ sites lie at each distance $d>\ell$ from $S$, the sum is bounded by
$4C\abs S^2\norm A\e^{vt}\sum_{d>\ell}\e^{-\mu d}$.
\end{proof}

\begin{figure*}[t]
\centering
\begin{adjustbox}{max width=\linewidth}
\begin{tenkz}[lattice={1x20}, frame=plane, physical=up, bonds=none]
  \tn{} & \tn{} & \tn{} & \tn{} & \tn{} & \tn{} & \tn{} & \tn{} & \tn{} & \tn{} & \tn{} & \tn{} & \tn{} & \tn{} & \tn{} & \tn{} & \tn{} & \tn{} & \tn{} & \tn{}
  \tnmark[form=enclosure, species=selected, label pos=s]{(1,5) .. (1,7)}{$R_0$}
  \tnmark[form=enclosure, species=secondary, label pos=n]{(1,3) .. (1,9)}{$R$}
  \tnmark[form=enclosure, species=collar, label pos=s]{(1,1) .. (1,2)}{$Q$}
  \tnmark[form=enclosure, species=collar, label pos=s]{(1,10) .. (1,12)}{$Q$}
  \tnmark[form=enclosure, species=complement, label pos=n]{(1,13) .. (1,20)}{$C$}
  \tnmark[form=enclosure, species=selected, label pos=s]{(1,15) .. (1,20)}{$C_0$}
  \tnmark[form=bracket, species=secondary, label pos=s]{(1,3) .. (1,4)}{$w$}
  \tnmark[form=bracket, species=secondary, label pos=s]{(1,8) .. (1,9)}{$w$}
  \tnmark[form=bracket, species=complement, label pos=s]{(1,13) .. (1,14)}{$w$}
  \tnmark[form=enclosure, species=neutral, tint, label pos=n]{(1,1) .. (1,3)}{$S$}
  \tnmark[form=enclosure, species=neutral, tint, label pos=n]{(1,9) .. (1,14)}{$S=N_\ell(Q\cup Z_R)$}
\end{tenkz}

\end{adjustbox}
\caption{The reset geometry, drawn in one replica of $\Psi=\Omega\otimes\Omega$: the cores $R_0\subset R$ and $C_0\subset C$, the buffer $Q=\Lambda\setminus(R\cup C)$, the margins of width $w$, the set $Z_R$ near the cut edges of $R$, and the support $S=N_\ell(Q\cup Z_R)$ of the unitary, which stays clear of both cores.}\label{fig:chain-reset}
\end{figure*}

\begin{lemma}[Information reset]\label{lem:reset}
For every $p>0$ there is $C_p<\infty$, depending only on $p,q,R,J,\Delta$,
with the following property for all sufficiently large $L$. Let
$R,C\subseteq\Lam$ be disjoint, $R_0\subseteq R$, $C_0\subseteq C$, and
suppose $I_\Om(R:C)\le b$. If
\[
\begin{split}
 &\dist(R_0,\Lam\setminus R)>w,\qquad \dist(C_0,\Lam\setminus C)>w,\\
 &w=C_p(1+b+\log L),
\end{split}
\]
then $I_\Om(R_0:C_0)\le L^{-p}$.
\end{lemma}
\begin{proof}
If either core is empty there is nothing to prove. Let $\Psi=\Om\otimes\Om$ on
two replicas, $F_R$ be the unitary exchanging the two replicas at the sites
of $R$, $\Phi=F_R\Psi$, and $Q=\Lam\setminus(R\cup C)$.

\emph{Overlap.} Apply Lemma~\ref{lem:splitting} to $(T,U,E)=(R,Q,C)$ to obtain an
isometry $V:\cH_Q\to\cB_R\otimes\cB_C$ and unit vectors $a\in\cH_R\otimes\cB_R$,
$a'\in\cH_C\otimes\cB_C$ satisfying
$\abs{\langle a\otimes a',\widetilde\Om\rangle}\ge\e^{-b/2}$, where
$\widetilde\Om=(\id\otimes V)\Om$. The largest Schmidt probability (marginal eigenvalue) of
$\widetilde\Om$ across the bipartition $R\cB_R\,|\,C\cB_C$ is therefore at least $\e^{-b}$,
so $\Tr\rho_{\widetilde\Om,R\cB_R}^2\ge\e^{-2b}$. On the two
replicas of $Q$, define the Hermitian contraction
$W=(V^\dagger\otimes V^\dagger)F_{\cB_R}(V\otimes V)$, where $F_{\cB_R}$ swaps
the two copies of $\cB_R$. Since $V$ acts outside $R$, the replica-swap identity
$\langle\xi^{\otimes2},F_{A}\xi^{\otimes2}\rangle=\Tr\rho_{\xi,A}^2$ gives
\begin{equation}\label{eq:zoverlap}
\begin{split}
 &z:=\langle\Phi,W\Psi\rangle=\langle\widetilde\Om^{\otimes2},F_RF_{\cB_R}\widetilde\Om^{\otimes2}\rangle\\
 &\qquad=\Tr\rho_{\widetilde\Om,R\cB_R}^2\ \ge\ \e^{-2b}.
\end{split}
\end{equation}

\emph{Two-sided Gaussian filtering.} Let $\mathbb H=H\otimes\id+\id\otimes H$
and $\mathbb H'=F_R\mathbb HF_R$. Both are range-$R$ Hamiltonians on the
doubled chain with uniformly bounded interaction sums, unique ground
vectors $\Psi$ and $\Phi$, common ground energy $\mathcal E=2E_0$, and spectral gap at
least $\Delta$. Their difference $\mathbb H'-\mathbb H$ is a sum of at most $N_R\abs{\partial R}$
terms of norm at most $2J$, namely those terms of $\mathbb H$ whose support meets both $R$
and $\Lam\setminus R$ (terms supported entirely inside $R$ are merely exchanged between the
two identical replicas); let $Z_R$ be the union of their supports, which lies
within distance $R$ of the cut edges of $R$. Put
$B_*=1+b+\log L$, $r_*=2p+10$, and $h=AB_*$ with $A\Delta^2/2\ge r_*+30$, and define
\[
\begin{split}
 &g_h(t)=(2\pi h)^{-1/2}\e^{-t^2/(2h)},\\
 &M=\int_{\R}g_h(t)\,\e^{it\mathbb H'}W\e^{-it\mathbb H}\,dt .
\end{split}
\]
The operator $M$ is a contraction, and by the spectral theorem
$M\Psi=\e^{-h(\mathbb H'-\mathcal E)^2/2}W\Psi$ and $M^\dagger\Phi=\e^{-h(\mathbb H-\mathcal E)^2/2}W^\dagger\Phi$.
Because $\mathbb H'$ and $\mathbb H$ have gap at least $\Delta$ above $\mathcal E$,
\begin{equation}\label{eq:twosided}
 \norm{M\Psi-z\Phi},\ \norm{M^\dagger\Phi-z\Psi}\le\e^{-h\Delta^2/2}\le\e^{-10b}L^{-r_*}.
\end{equation}
Truncating the time integral to $\abs t\le T=T_0B_*$ incurs an error of at most
$2\e^{-T_0^2B_*/(2A)}$, which is likewise at most $\e^{-10b}L^{-r_*}$ once
$T_0^2/(2A)\ge r_*+30$.

\emph{Localisation.} Factor the integrand as
$\e^{it\mathbb H'}W\e^{-it\mathbb H}=(\e^{it\mathbb H'}W\e^{-it\mathbb H'})\,V(t)$
with $V(t)=\e^{it\mathbb H'}\e^{-it\mathbb H}$, whose left generator
$G(t)=\e^{it\mathbb H'}(\mathbb H'-\mathbb H)\e^{-it\mathbb H'}$ is a sum of
the boundedly many terms supported in $Z_R$. Applying Lemma~\ref{lem:LR} to $W$ (supported in $Q$, of norm at most $1$)
and to each term of $G(t)$ yields Hermitian localised approximants
$W_\ell(t)=\mathcal P_{N_\ell(Q)}\bigl(\e^{it\mathbb H'}W\e^{-it\mathbb H'}\bigr)$
and $G_\ell(t)=\mathcal P_{N_\ell(Z_R)}G(t)$ satisfying
\[
 \norm{\e^{it\mathbb H'}W\e^{-it\mathbb H'}-W_\ell(t)}+\norm{G(t)-G_\ell(t)}\le CL^2\,\e^{v\abs t-\mu\ell},
\]
where the factor $L^2$ bounds $\abs Q^2$ and $\abs{\partial R}$. Let
$V_\ell(t)$ be the unitary evolution generated by $G_\ell(t)$, which is supported in
$N_\ell(Z_R)$; Duhamel's formula gives $\norm{V(t)-V_\ell(t)}\le CL^2\abs t\e^{v\abs t-\mu\ell}$.
Consequently, $W_\ell(t)V_\ell(t)$ is a contraction supported in
$S:=N_\ell(Q\cup Z_R)$ that approximates the integrand uniformly for
$\abs t\le T$ with error at most $CL^2(1+T)\e^{vT-\mu\ell}$. Choose
$\ell=\lceil K_{\mathrm{loc}}B_*\rceil$ with $K_{\mathrm{loc}}$ fixed large
enough (after $T_0$) that this error is at most $\e^{-10b}L^{-r_*}$,
using $\log L\le B_*$ and $\log(1+T)\le C+B_*$. Averaging the localised
integrands against the truncated Gaussian yields a contraction
$M_{\mathrm{loc}}$ supported in $S$ with
\begin{equation}\label{eq:loctwosided}
 \norm{M_{\mathrm{loc}}\Psi-z\Phi},\ \norm{M_{\mathrm{loc}}^\dagger\Phi-z\Psi}\le C\e^{-10b}L^{-r_*}.
\end{equation}

\emph{Polar decomposition.} Rescale by setting $D=M_{\mathrm{loc}}/z$. Since $z\ge\e^{-2b}$, we have $\norm D\le\e^{2b}$ and
$\norm{D\Psi-\Phi},\norm{D^\dagger\Phi-\Psi}\le C\e^{-8b}L^{-r_*}$. Writing
$(D^\dagger D-\id)\Psi=D^\dagger(D\Psi-\Phi)+(D^\dagger\Phi-\Psi)$ then gives
$\norm{(D^\dagger D-\id)\Psi}\le C\e^{-6b}L^{-r_*}$. Because $D$ acts on the finite-dimensional space $\cH_S$, its polar partial
isometry extends to a unitary $U$ on
$\cH_S$ tensored with the
identity outside $S$. Using the scalar inequality $(\sqrt x-1)^2\le(x-1)^2$ for $x\ge0$, we obtain
$\norm{(U-D)\Psi}=\norm{(\id-\abs D)\Psi}\le\norm{(\id-D^\dagger D)\Psi}$, and therefore
\begin{equation}\label{eq:Uerror}
 \norm{U\Psi-\Phi}\le C\e^{-6b}L^{-r_*}\le CL^{-r_*}.
\end{equation}

\emph{Support and conclusion.} By hypothesis, every site of $Q$ lies at distance greater than $w$ from
both $R_0$ and $C_0$. Any interaction term whose support meets both $R$ and $\Lam\setminus R$
contains a site $y\notin R$ (at distance greater than $w$ from $R_0$) and a site
$x\in R\subseteq\Lam\setminus C$ (at distance greater than $w$ from $C_0$); since its support
has diameter at most $R$, every site of $Z_R$ lies at distance greater than $w-R$ from
both cores. Choosing $C_p>K_{\mathrm{loc}}+R+3$ ensures that $S=N_\ell(Q\cup Z_R)$ is
disjoint from $R_0\cup C_0$ in both replicas. Thus $U$ leaves the
two-replica marginal on $R_0C_0$ unchanged. In $\Phi=F_R\Psi$, the marginal on
$R_0C_0$ of the \emph{first} replica is $\rho_{R_0}\otimes\rho_{C_0}$,
because $R_0$ is swapped in from the second replica whereas $C_0$ remains in the first. By
\eqref{eq:Uerror} and the contractivity of the partial trace, the ground-state
marginal $\rho_{R_0C_0}$ lies within trace distance $CL^{-r_*}$ of
$\rho_{R_0}\otimes\rho_{C_0}$. Applying the Alicki--Fannes--Winter continuity bound
$\abs{S(A|B)_\rho-S(A|B)_\sigma}\le C\sqrt\delta\log(\e\dim\cH_A)$ with
$\delta=\frac12\norm{\rho-\sigma}_1$ \cite{Winter2016} to
$I(R_0:C_0)=S(R_0)-S(R_0|C_0)$, and using $\log\dim\cH_{R_0}\le L\log q$, gives
$I_\Om(R_0:C_0)\le CL^{1-r_*/2}=CL^{-p-4}\le L^{-p}$ for all sufficiently large $L$.
\end{proof}

\begin{figure*}[t]
\centering
\begin{adjustbox}{max width=\linewidth}
\(\displaystyle \Phi=F_R\Psi\;=\;
  \begin{tenkz}[rows={wire,wire}, cols=7, bonds=none]
    \tn[at=(2,1), skin=box, wide=3, name=o1,
        ports={90@1:physical:$C$, 90@2:physical:$Q$, 90@3:physical}]{\Omega}
    \tn[at=(2,5), skin=box, wide=3, name=o2,
        ports={90@1:physical, 90@2:physical:$Q$, 90@3:physical:$C$}]{\Omega}
    \tn[at=(1,3), skin=box, wide=3, name=f,
        ports={270@1:physical, 270@3:physical, 90@1:physical:$R$, 90@3:physical:$R$}]{F_R}
    \tnwire{o1.90@3}{f.270@1}
    \tnwire{o2.90@1}{f.270@3}
  \end{tenkz}
  \;=\;
  \begin{tenkz}[rows={wire,wire}, cols=7, bonds=none]
    \tn[at=(2,1), skin=box, wide=3, name=o1,
        ports={90@1:physical:$C$, 90@2:physical:$Q$, 90@3:physical}]{\Omega}
    \tn[at=(2,5), skin=box, wide=3, name=o2,
        ports={90@1:physical, 90@2:physical:$Q$, 90@3:physical:$C$}]{\Omega}
    \tn[at=(1,3), skin=none, name=r1, ports={270:physical, 90:physical:$R$}]{}
    \tn[at=(1,5), skin=none, name=r2, ports={270:physical, 90:physical:$R$}]{}
    \tnwire[name=a]{o1.90@3}{r2.270}
    \tnwire[name=b, cross={under at crossing of a and b}]{o2.90@1}{r1.270}
  \end{tenkz}\)

\end{adjustbox}
\caption{The swap $F_R$ exchanges the two copies of $\Omega$ on $R$ only (the first copy is drawn mirrored so that the two $R$ legs are adjacent). In $\Phi=F_R\Psi$ the core $R_0$ of the first copy comes from the second copy while $C_0$ does not, so the first copy's marginal on $R_0C_0$ is $\rho_{R_0}\otimes\rho_{C_0}$.}\label{fig:replica-swap}
\end{figure*}

\begin{figure*}[t]
\centering
\begin{adjustbox}{max width=\linewidth}
\(\displaystyle M=\int_{\mathbb R}g_h(t)\,dt\;
  \begin{quantikz}[row sep={0.55cm,between origins}]
    \lstick{$R^{(1)}$} & \gate[6]{e^{-it\mathbb H}} & & \gate[6]{e^{it\mathbb H'}} & \\
    \lstick{$R^{(2)}$} & & & & \\
    \lstick{$Q^{(1)}$} & & \gate[2]{W} & & \\
    \lstick{$Q^{(2)}$} & & & & \\
    \lstick{$C^{(1)}$} & & & & \\
    \lstick{$C^{(2)}$} & & & &
  \end{quantikz}\)

\end{adjustbox}
\caption{The two-sided Gaussian filter: evolve with the doubled Hamiltonian, apply $W$ on the two copies of the buffer, evolve back with the swapped Hamiltonian, and average against a Gaussian in time. It sends $W\Psi$ to a multiple of $\Phi$ and $W^\dagger\Phi$ to a multiple of $\Psi$, up to $e^{-h\Delta^2/2}$.}\label{fig:qc-gaussian-filter}
\end{figure*}

\subsection{The one-dimensional separation lemma}\label{sec:logsep2}

\begin{inwords}
Combining the one-dimensional area law with the information reset yields
polynomial decoupling at logarithmic separation. By Theorem~\ref{thm:A}, any
interval $R$ (or union of $k_0$ intervals) has entanglement entropy at most
$k_0C_{\mathrm{int}}$, so its mutual information with any disjoint region $C$
is bounded by the constant $2k_0C_{\mathrm{int}}$. Feeding this uniform bound
into Lemma~\ref{lem:reset} shows that any union $T$ of at most $k_0$ intervals
and any set $E$ separated from $T$ by a distance of order $\log L$ satisfy
$I_\Om(T:E)\le L^{-p}$. Because the one-dimensional area law bounds interval
entropies uniformly at every scale, this single step establishes logarithmic
decoupling for intervals of arbitrary length, from a few sites up to the
entire chain. This replaces the entire Section~3 of the two-dimensional
paper (coarse collar, cell conditional information, angular induction).
\end{inwords}

\precisely
\begin{lemma}[Information at logarithmic separation]\label{lem:logsep}
For every $p>0$ and $k_0\ge1$ there is $D_p=D_p(q,R,J,\Delta,k_0)$ such that for
all sufficiently large $L$, every $T\subseteq\Lam$ that is a union of at most
$k_0$ intervals, and every $E\subseteq\Lam$ with $\dist(T,E)\ge D_p\log L$,
\[
 I_\Om(T:E)\le L^{-p}.
\]
\end{lemma}
\begin{proof}
Let $b=4k_0C_{\mathrm{int}}$, let $w=C_p(1+b+\log L)$ be the margin width from Lemma~\ref{lem:reset},
and choose $D_p$ so that $D_p\log L\ge 3w+3$. Put $R=T^{+(w+1)}$ (a union of at most $k_0$
intervals), $C=\Lam\setminus T^{+(2w+2)}$, $R_0=T$, and $C_0=E$. Then
$\dist(R_0,\Lam\setminus R)\ge w+1$ and $\dist(C_0,\Lam\setminus C)\ge w+1$. Moreover,
$I(R:C)\le 2S(R)\le 2k_0C_{\mathrm{int}}\le b$ by Theorem~\ref{thm:A},
using subadditivity across the at most $k_0$ intervals of $R$ together with
$I(R:C)=S(R)+S(C)-S(RC)\le 2S(R)$ from the Araki--Lieb inequality.
Applying Lemma~\ref{lem:reset} gives $I_\Om(T:E)\le L^{-p}$.
\end{proof}

\begin{figure*}[t]
\centering
\begin{adjustbox}{max width=\linewidth}
\begin{tenkz}[lattice={1x18}, frame=plane, physical=up, bonds=none]
  \tn{} & \tn{} & \tn{} & \tn{} & \tn{} & \tn{} & \tn{} & \tn{} & \tn{} & \tn{} & \tn{} & \tn{} & \tn{} & \tn{} & \tn{} & \tn{} & \tn{} & \tn{}
  \tnmark[form=enclosure, species=selected, label pos=s]{(1,8) .. (1,10)}{$T$}
  \tnmark[form=enclosure, species=secondary, label pos=n]{(1,6) .. (1,12)}{$R=T^{+(w+1)}$}
  \tnmark[form=enclosure, species=complement, label pos=n]{(1,1) .. (1,3)}{$C$}
  \tnmark[form=enclosure, species=complement, label pos=n]{(1,15) .. (1,18)}{$C$}
  \tnmark[form=enclosure, species=collar, label pos=s]{(1,1) .. (1,2)}{$E$}
  \tnmark[form=enclosure, species=collar, label pos=s]{(1,17) .. (1,18)}{$E$}
  \tnmark[form=bracket, species=collar, label pos=s]{(1,3) .. (1,7)}{$D_p\log L$}
\end{tenkz}

\end{adjustbox}
\caption{The separation lemma: an interval $T$, its enlargement $R$, the far region $C$ and the set $E\subseteq C$ at distance $D_p\log L$ from $T$. The area law bounds $I(R:C)$ by a constant and the reset makes $I(T:E)\le L^{-p}$.}\label{fig:chain-logsep}
\end{figure*}

\begin{remark}[A route without Lieb--Robinson bounds]\label{rem:LRfree}
The information-reset lemma is the only place in Part~\ref{part:mps} where a
Lieb--Robinson bound is invoked. It can be replaced by the algebraic amplification
argument of \cite[Section~10]{AL}, which uses exact positive quasi-local
constraints with stretched-exponential tails $\e^{-cl^\alpha}$
($\alpha=p/(p+1)$), random products of the associated contractions, and a
polar decomposition, without requiring a propagation bound for the unitary
dynamics. With the area law as input, that argument yields
$I(T:E)\le L^{-M}$ whenever $E$ lies beyond distance $C(\log L)^{1/\alpha}$ of an
interval $T$. Running the protocol below on blocks of that width produces
a matrix product state of bond dimension $\exp(O((\log L)^{1/\alpha}))$, which is
quasi-polynomial for every fixed $p$ and approaches polynomial only in the large-$p$ limit.
Achieving strictly polynomial bond dimension through the protocol route relies on the
logarithmic separation provided by Lemma~\ref{lem:logsep}.
\end{remark}
